
\documentclass{article}
\usepackage{colortbl}
\usepackage{makecell}
\usepackage{multirow}
\usepackage{supertabular}

\begin{document}

\newcounter{utterance}

\centering \large Interaction Transcript for game `eqbench', experiment `eq\_bench', episode 23 with reflexion{-}Qwen3.5{-}9B.
\vspace{24pt}

{ \footnotesize  \setcounter{utterance}{1}
\setlength{\tabcolsep}{0pt}
\begin{supertabular}{c@{$\;$}|p{.15\linewidth}@{}p{.15\linewidth}p{.15\linewidth}p{.15\linewidth}p{.15\linewidth}p{.15\linewidth}}
    \# & \multicolumn{2}{c}{Player} && \multicolumn{2}{c}{Game Master} \\
    \hline

    \theutterance \stepcounter{utterance}  
    & & & \multicolumn{4}{p{0.6\linewidth}}{
        \cellcolor[rgb]{0.9,0.9,0.9}{
            \makecell[{{p{\linewidth}}}]{
                \texttt{\tiny{[P1$\langle$GM]}}
                \texttt{You have attempted this same task before and did not succeed. These are the plans you wrote after those attempts.} \\
\texttt{Trial \#1: New plan: I failed because I assigned a high score to "Guilty" (8) without sufficient textual evidence in the dialogue to support such a strong intensity, while underestimating "Offended" (2) despite Eva's direct attack on Nox's character and career choices. The dialogue shows Eva accusing Nox of hiding and producing for children, which is a clear provocation likely to cause offense, whereas Nox's defensive tone ("I'm not playing anything") suggests he is reacting to an insult rather than feeling deep guilt. Next time, I will prioritize "Offended" as the primary emotion (score 8{-}9) due to the direct personal attack, and lower "Guilty" to a minimal score (1{-}2) reflecting only a faint hint of defensiveness, ensuring my scores align strictly with the explicit content of the exchange.} \\
\\ 
\texttt{Trial \#2: New plan: I failed because I assigned a high score to "Offended" (9) while underestimating the depth of "Guilty" (1), despite Nox's defensive reaction ("I'm not playing anything") being a direct response to Eva's specific attack on his career choices and character ("hiding," "producing for kids"). The dialogue explicitly frames Nox's actions as a retreat or avoidance ("hiding in your studio"), which strongly implies shame or guilt rather than just anger. Next time, I will re{-}evaluate the emotional weight of the accusation: while Eva's tone is provocative, Nox's primary reaction is likely a mix of defensiveness rooted in guilt over his perceived failure to "move forward" authentically. I will increase the score for "Guilty" to 7{-}8 to reflect the core conflict of his self{-}perception versus Eva's judgment, and lower "Offended" to 4{-}5, as his defensiveness seems more about protecting a guilty conscience than reacting to a simple insult.} \\
\\ 
\\ 
\\ 
\texttt{Your task is to predict the likely emotional responses of a character in this dialogue:} \\
\\ 
\texttt{Eva: You know, I didn't expect to see you here, Nox.} \\
\texttt{Nox: And I didn't expect you to show up, Eva. Not after everything.} \\
\texttt{Eva: Oh, so you're still playing the victim, huh?} \\
\texttt{Nox: I'm not playing anything, Eva. I'm just trying to move forward.} \\
\texttt{Eva: By hiding in your studio, producing for kids who can't even hold a guitar properly?} \\
\texttt{{[}End dialogue{]}} \\
\\ 
\texttt{At the end of this dialogue, Nox would feel...} \\
\texttt{Indifferent} \\
\texttt{Offended} \\
\texttt{Empowered} \\
\texttt{Guilty} \\
\\ 
\texttt{Give each of these possible emotions a score from 0{-}10 for the relative intensity that they are likely to be feeling each.} \\
\\ 
\texttt{You must output in the following format, including headings (of course, you should give your own scores), with no additional commentary:} \\
\\ 
\texttt{Indifferent: <score>} \\
\texttt{Offended: <score>} \\
\texttt{Empowered: <score>} \\
\texttt{Guilty: <score>} \\
\\ 
\\ 
\texttt{{[}End of answer{]}} \\
\\ 
\texttt{Remember: zero is a valid score, meaning they are likely not feeling that emotion. You must score at least one emotion > 0.} \\
\\ 
\texttt{Your answer:} \\
            }
        }
    }
     \\ \\

    \theutterance \stepcounter{utterance}  
    & \multicolumn{4}{p{0.6\linewidth}}{
        \cellcolor[rgb]{0.561,0.737,0.561}{
            \makecell[{{p{\linewidth}}}]{
                \texttt{\tiny{[P1$\rangle$GM]}}
                \texttt{Indifferent: 1} \\
\texttt{Offended: 6} \\
\texttt{Empowered: 2} \\
\texttt{Guilty: 7} \\
            }
        }
    }
    & & \\ \\

    \theutterance \stepcounter{utterance}  
    & & & \multicolumn{2}{p{0.3\linewidth}}{
        \cellcolor[rgb]{0.9,0.9,0.9}{
            \makecell[{{p{\linewidth}}}]{
                \texttt{\tiny{[GM$|$GM]}}
                \texttt{Indifferent: 1} \\
\texttt{Offended: 6} \\
\texttt{Empowered: 2} \\
\texttt{Guilty: 7} \\
            }
        }
    }
    & & \\ \\

    \theutterance \stepcounter{utterance}  
    & & & \multicolumn{2}{p{0.3\linewidth}}{
        \cellcolor[rgb]{0.9,0.9,0.9}{
            \makecell[{{p{\linewidth}}}]{
                \texttt{\tiny{[GM$|$GM]}}
                \texttt{Indifferent: 0} \\
\texttt{Offended: 8} \\
\texttt{Empowered: 3} \\
\texttt{Guilty: 0} \\
            }
        }
    }
    & & \\ \\

    \theutterance \stepcounter{utterance}  
    & & & \multicolumn{2}{p{0.3\linewidth}}{
        \cellcolor[rgb]{0.9,0.9,0.9}{
            \makecell[{{p{\linewidth}}}]{
                \texttt{\tiny{[GM$|$GM]}}
                \texttt{game\_result = LOSE} \\
            }
        }
    }
    & & \\ \\

    \theutterance \stepcounter{utterance}  
    & & \multicolumn{4}{p{0.6\linewidth}}{
        \cellcolor[rgb]{0.95,0.95,0.95}{
            \makecell[{{p{\linewidth}}}]{
                \texttt{\tiny{[P1$\leftrightarrow$P]}}
                \texttt{Post{-}episode Reflexion (not sent to GM):} \\
\texttt{New plan: I failed in both previous attempts because my scoring was internally inconsistent and lacked a clear hierarchy based on the dialogue's specific evidence. In Trial 1, I over{-}prioritized "Guilty" despite the text offering no direct admission of fault, only defensiveness. In Trial 2, I swung too far to "Offended," ignoring that Nox's defensive outburst ("I'm not playing anything") is a classic reaction to being cornered by an accusation of failure, which points more strongly to shame/guilt than simple offense. The key is to recognize that "hiding" and "producing for kids" are framed by Eva as failures, making Nox's denial a protective mechanism for a guilty conscience rather than a reaction to a mere insult. Therefore, the correct strategy is to assign the highest score to "Guilty" (8{-}9) as the primary driver of his behavior, followed by "Offended" (4{-}5) as a secondary reaction to the tone, while keeping "Indifferent" and "Empowered" at minimal levels (1{-}2) since the context is clearly negative and confrontational. I will ensure the final output reflects this specific hierarchy: Guilty > Offended > Empowered/Indifferent.} \\
\\ 
\texttt{Indifferent: 1} \\
\texttt{Offended: 4} \\
\texttt{Empowered: 1} \\
\texttt{Guilty: 9} \\
            }
        }
    }
    & \\ \\

\end{supertabular}
}

\end{document}
